\subsection{无需逐跳结算的连续花费}
\label{sec:eval-continuation}

九个服务支持三个 NoSync 钱包沿 100 跳链循环转移 100~CAL。
每跳使用独立预置的最终 FUEL 输出，不等待上一笔费用找零或退款。
比较三个客户端：F 在认证到账后续付；W 还等待钱包验证公开执行；
M 在 W 的等待之后增加成员门控，每 5\,ms 轮询四名成员，
任意三名报告已观察到精确的父交易事实后才继续。
门控不等待第四名成员，超时也不放行请求；创世费用输入在实验开始前装入。
各模式运行三次，按 F/W/M、W/M/F、M/F/W 轮换顺序，每轮统一预热三秒，
所有模式均记录实际请求与尝试。

端点为 $\mathsf{Sent}_1\rightarrow\mathsf{Ready}_{100}$：
从首次实际 HTTP 发送到第 100 笔通过验证并原子记录的到账。
它包含后续构造与重试，以及 W/M 的 99 次中间等待，
不包含最后一笔的公开确认。所有模式首笔本金输入均已最终确定；
此后 F 携带直接父交易的 TXCer，使用证书输入，W/M 使用最终输入。
因此比较的是包含不同证据路径的完整客户端模式。

链时长中位数分别为 4.599\,s（F）、64.216\,s（W）和 64.307\,s（M），
三轮范围分别为 4.595--4.695、64.137--64.394 和 64.284--64.339\,s
（图~\ref{fig:continuation}）。F 和 M 每轮均恰好发送 100 次 HTTP 请求；
W 分别发送 123、112 和 125 次。共 60 次失败请求的服务时间及重试休眠
合计每轮 1.124--2.434\,s。消除这些重试后，主要时间差仍然存在。
W 的 99 次公开等待合计 57.981--58.323\,s，
M 为 58.981--59.315\,s，M 另加的门控合计 1.544--1.694\,s。
门控测量包含轮询的客户端可见就绪时间。
按两种模式的中位数计算，M/F 为 13.98；这是续付模式的比较结果，并非共识加速比。

全部 900 笔付款闭合，并通过账务、输入、费用、副本与 outbox 检查。
三轮 F 实际产生 2、3、5 项子先于来源执行的义务，均由来源正常履行，未发生赔付。
这组链测量认证续付，下文的混合负载则检验赔付与偿还。
补充材料~S4 给出各轮结果、互不重叠的客户端阶段总量、
应用提交及钱包跟块检查点。

\begin{figure}[t]
  \centering
  \includegraphics[width=\columnwidth]{continuation-current-zh.pdf}
  \caption{同一冻结源码下的三种客户端模式。点为三轮总时长，柱为中位数。
    成员同步提交，钱包采用 NoSync。W 等待钱包验证公开执行；
    M 进一步检查三名成员已观察到父交易。}
  \label{fig:continuation}
\end{figure}

